\chapter{Conclusão e Trabalho Futuro}
\label{cap:conclusao}

\section{Requisitos Cumpridos}

A Tabela~\ref{tab:rastreabilidade} relaciona cada requisito definido na Secção~2.3 com a sua
implementação.

\begin{longtable}{p{0.09\textwidth}p{0.13\textwidth}p{0.66\textwidth}}
\label{tab:rastreabilidade} \\
\toprule
\textbf{Req.} & \textbf{Estado} & \textbf{Implementação} \\
\midrule
\endfirsthead
\toprule
\textbf{Req.} & \textbf{Estado} & \textbf{Implementação} \\
\midrule
\endhead
\bottomrule
\endfoot
RF01 & Cumprido & Registo de clientes e login (Auth); gestores criados pelo administrador \\
RF02 & Cumprido & Recuperação de palavra-passe com ligação por email (Auth e Notifications) \\
RF03 & Cumprido & Pesquisa por texto (nome, cozinha, cidade), cozinha, data e número de pessoas;
a hora escolhe-se na disponibilidade (Restaurants) \\
RF04 & Cumprido & Detalhes do restaurante e horas disponíveis, incluindo as esgotadas
(Restaurants e Reservations) \\
RF05 & Cumprido & Reserva com data, hora e número de pessoas, sem exceder as mesas \\
RF06 & Cumprido & Consultar (próximas e histórico), alterar e cancelar reservas \\
RF07 & Cumprido & Entrar e sair da fila, posição, total de grupos e tempo estimado \\
RF08 & Cumprido & Emails de reserva confirmada, alterada e cancelada e de mesa disponível para o
grupo da fila (Notifications) \\
RF09 & Cumprido & Reservas do restaurante, reservas manuais, chegada com sugestão de mesa e não
comparência \\
RF10 & Cumprido & Fila do restaurante, chamar o próximo grupo que cabe e atribuição automática de
mesas livres \\
RF11 & Cumprido & Mesas, capacidades, turnos, horários e tempo médio de refeição \\
RF12 & Parcial & Ocupação por dia e tempo médio de refeição real (Restaurants); taxa de faltas
nos relatórios do administrador. Faltam estatísticas dos tempos de espera da fila \\
RF13 & Cumprido & Criar e editar restaurantes; remover é feito por desativação, para manter o
histórico \\
RF14 & Cumprido & Lista de utilizadores com filtros, bloqueio e desbloqueio; os clientes que
atingem o limite de faltas ficam marcados para revisão \\
RF15 & Cumprido & Limite de faltas (Auth) e tempo de refeição, margem e tolerância por defeito
(Restaurants) \\
RF16 & Cumprido & Reservas e taxa de faltas por restaurante ou por região \\
RF17 & Cumprido & Alterar palavra-passe e fotografia de perfil \\
\midrule
RNF01 & Por medir & Estimativa atualizada por eventos, em menos de um segundo após cada alteração;
não foram feitos testes de carga para medir o percentil 95 \\
RNF02 & Parcial & JWT~\cite{rfc7519} validado no gateway e em cada serviço, BCrypt~\cite{bcrypt}, controlo de acesso por perfil.
O ambiente local usa HTTP; o HTTPS fica a cargo do ponto de entrada em produção \\
RNF03 & Cumprido & Só o gateway publica uma porta; serviços, base de dados e RabbitMQ~\cite{rabbitmq} na rede
interna \\
RNF04 & Cumprido & Comunicação assíncrona; demonstrado com o serviço de notificações parado
(Tabela~\ref{tab:falhas}) \\
RNF05 & Cumprido & Os microsserviços devem ser executados em containers Docker~\cite{docker} e não guardar estado em memória: os dados ficam na respetiva base de dados PostgreSQL~\cite{postgresql} e as mensagens no RabbitMQ~\cite{rabbitmq}.\\
RNF06 & Cumprido & Um serviço e uma base de dados por área de negócio, \textit{migrations}
versionadas e OpenAPI~\cite{openapi} de cada serviço \\
RNF07 & Protótipo & Abordagem \textit{mobile first} definida no protótipo de alta fidelidade; a
aplicação cliente não fez parte da implementação \\
RNF08 & Cumprido & 114 testes unitários e de integração e duas \textit{collections} do Postman~\cite{postman}
(Capítulo~\ref{cap:testes}) \\
RNF09 & Cumprido & Consentimento no registo e eliminação de conta com anonimização em todos os
serviços \\
\caption{Rastreabilidade dos requisitos}
\end{longtable}

\section{Conclusão}

O GrabATable foi implementado como um conjunto de cinco microsserviços independentes atrás de um
API Gateway, cada um com a sua base de dados, comunicando por eventos através do RabbitMQ~\cite{rabbitmq}. A
principal inovação identificada na análise, a gestão conjunta das reservas e da fila de espera, foi
concretizada sem transações distribuídas: cada decisão crítica (não exceder as mesas, contar uma
falta uma só vez, não chamar o mesmo grupo duas vezes) é tomada dentro de uma única base de dados,
e a informação de que a fila de espera precisa chega-lhe por eventos.

A arquitetura foi validada não só no caminho feliz, mas também perante falhas: com qualquer
serviço, a base de dados ou o \textit{broker} em baixo, o sistema responde de forma previsível e
explica à aplicação o que falhou, e o rastreio distribuído permite seguir cada pedido através de
todos os serviços. O comportamento está coberto por 114 testes automáticos e por testes de sistema
que percorrem os fluxos completos e os cenários de falha.

\section{Trabalho Futuro}

\begin{itemize}
  \item \textbf{\textit{Transactional outbox}:} garantir que nenhum evento se perde quando o RabbitMQ~\cite{rabbitmq}
        está em baixo;
  \item \textbf{Aplicação cliente:} implementar a aplicação a partir do protótipo, com notificações
        em tempo real (por exemplo com SignalR) para a posição na fila;
  \item \textbf{Estatísticas da fila:} tempos de espera reais face aos estimados, para completar o
        RF12 e afinar a estimativa;
  \item \textbf{Desempenho:} testes de carga para verificar o RNF01 e calibrar o
        \textit{HorizontalPodAutoscaler};
  \item \textbf{Segurança em produção:} HTTPS no ponto de entrada e gestão de segredos num cofre;
  \item \textbf{Publicação numa \textit{cloud}:} publicar o projeta na cloud.
\end{itemize}
